\documentclass[11pt]{article}

\usepackage[margin=1in]{geometry}
\usepackage{amsmath,amssymb}
\usepackage{booktabs}
\usepackage{array}

\title{Why the Two Flexoelectric Potential Formulas Look the Same}
\author{}
\date{}

\begin{document}
\maketitle

\section{Short Answer}

The two potential expressions have the same visible shape because both come from
integrating a through-thickness electric field of the form
\[
  E_2(Y)=\hbox{constant}
  +\hbox{constant}\,
  \frac{\cosh(Y/\ell_b)}{\cosh(a/\ell_b)},
  \qquad a=\frac{H}{2}.
\]
After integration, both potentials must have the form
\[
  \phi(Y)
  =
  \hbox{linear term}
  +
  \hbox{sinh boundary-layer term}
  +
  \hbox{constant}.
\]
Therefore, the visual form of the final potential is not enough to identify the
mechanical assumption used in the derivation.  The difference is mainly in the
coefficients that define the boundary-layer length and amplitude.

\section{Same Shape, Different Voltage Reference}

The older expression may be written as
\begin{equation}
  \phi_{\mathrm{old}}(Y)
  =
  -A_0^{\mathrm{old}}Y
  +
  B_0^{\mathrm{old}}
  \frac{\sinh(Y/\ell_b)}{\cosh(a/\ell_b)}
  +
  \phi_0 .
  \label{eq:old-shape}
\end{equation}
Here \(\phi_0\) is arbitrary because electric potential is defined only up to a
constant.

The new generalized-traction expression was written with the bottom surface as
the voltage reference:
\[
  \phi_{\mathrm{new}}(-a)=0.
\]
It is
\begin{equation}
  \phi_{\mathrm{new}}(Y)
  =
  -A_0^{\mathrm{new}}(Y+a)
  +
  B_0^{\mathrm{new}}
  \left[
    \frac{\sinh(Y/\xi)}{\cosh(a/\xi)}
    +
    \tanh(a/\xi)
  \right].
  \label{eq:new-bottom-gauge}
\end{equation}
Expanding Eq.~\eqref{eq:new-bottom-gauge} gives
\begin{equation}
  \phi_{\mathrm{new}}(Y)
  =
  -A_0^{\mathrm{new}}Y
  +
  B_0^{\mathrm{new}}
  \frac{\sinh(Y/\xi)}{\cosh(a/\xi)}
  +
  \left[
    -A_0^{\mathrm{new}}a
    +
    B_0^{\mathrm{new}}\tanh(a/\xi)
  \right].
  \label{eq:new-expanded}
\end{equation}
Thus Eq.~\eqref{eq:new-bottom-gauge} is the same shape as
Eq.~\eqref{eq:old-shape}.  The extra \(\tanh(a/\xi)\) term only fixes the
constant so that \(\phi(-a)=0\).  It is not an additional physical field.

\section{Where the Shorthand Definitions Differ}

The two derivations use the same kind of basis function, but they do not
necessarily use the same coefficients.  The safest way to compare the formulas
is not to start from \(\ell_b\) and \(\xi\).  Instead, compare the coefficients
from which those lengths are built.

Define
\[
  C_L=\lambda+2\mu,\qquad
  C_T=\lambda,\qquad
  \beta=\frac{C_T}{C_L}.
\]
Then the bulk part of the electric field gives
\[
  A_0=\frac{\kappa}{\epsilon}\left(\mu_T-\mu_L\beta\right).
\]
This part is common.  The difference is in the shorthand used for the
boundary-layer part.

\[
\begin{array}{>{\raggedright\arraybackslash}p{0.20\linewidth}
                >{\raggedright\arraybackslash}p{0.36\linewidth}
                >{\raggedright\arraybackslash}p{0.36\linewidth}}
\toprule
Quantity & Older shorthand & New generalized-traction shorthand \\
\midrule
Primary coefficient &
\[
  H_L^*
  =
  C_L\ell^2
  -
  \frac{\epsilon_0\mu_L^2}
       {\epsilon(\epsilon-\epsilon_0)}
\]
&
\[
  G_L
  =
  H_L+\frac{\mu_L^2}{\epsilon}
\]
\\
Transverse coefficient &
\[
  H_T^*
  =
  C_T\ell^2
  -
  \frac{\epsilon_0\mu_L\mu_T}
       {\epsilon(\epsilon-\epsilon_0)}
\]
&
\[
  G_T
  =
  H_T+\frac{\mu_L\mu_T}{\epsilon}
\]
\\
Bare model substitution &
\[
  \begin{aligned}
  H_L^*&=C_L\ell^2,\\
  H_T^*&=C_T\ell^2
  \quad(\epsilon_0=0)
  \end{aligned}
\]
&
\[
  \begin{aligned}
  G_L&=C_L\ell^2+\frac{\mu_L^2}{\epsilon},\\
  G_T&=C_T\ell^2+\frac{\mu_L\mu_T}{\epsilon}
  \end{aligned}
\]
\\
\bottomrule
\end{array}
\]

The starred quantities \(H_L^*,H_T^*\) are the coefficients used in the older
reduced formula.  The quantities \(G_L,G_T\) are the effective coefficients that
appear in the higher-order stress after the electric field is eliminated from
the open-circuit condition \(D_2=0\):
\begin{equation}
  G_L=H_L+\frac{\mu_L^2}{\epsilon},
  \qquad
  G_T=H_T+\frac{\mu_L\mu_T}{\epsilon}.
  \label{eq:G}
\end{equation}
For the default bare generalized-traction model used by the plotting script,
\begin{equation}
  H_L=C_L\ell^2,\qquad H_T=C_T\ell^2,
\end{equation}
so
\begin{equation}
  G_L=C_L\ell^2+\frac{\mu_L^2}{\epsilon},
  \qquad
  G_T=C_T\ell^2+\frac{\mu_L\mu_T}{\epsilon}.
  \label{eq:G-bare}
\end{equation}
Therefore \(H_L^*\) and \(G_L\) are not just two symbols for the same definition.
In the bare case with \(\ell=0\),
\[
  H_L^*=0,
  \qquad
  G_L=\frac{\mu_L^2}{\epsilon}.
\]
So the older shorthand gives no real boundary-layer length from \(H_L^*\), while
the generalized-traction shorthand still gives a finite length if \(\mu_L\ne0\).

\section{How the Different Coefficients Propagate}

Once \(H_L^*,H_T^*\) or \(G_L,G_T\) are chosen, the rest of the shorthand follows.
This is where the formulas look deceptively similar:

\[
\begin{array}{>{\raggedright\arraybackslash}p{0.20\linewidth}
                >{\raggedright\arraybackslash}p{0.36\linewidth}
                >{\raggedright\arraybackslash}p{0.36\linewidth}}
\toprule
Quantity & Older expression & New generalized-traction expression \\
\midrule
Boundary length &
\[
  \ell_b=\sqrt{\frac{H_L^*}{C_L}}
\]
&
\[
  \xi=\sqrt{\frac{G_L}{C_L}}
\]
\\
Boundary mismatch &
\[
  \eta_{\mathrm{old}}
  =
  \frac{H_T^*}{H_L^*}
  -
  \frac{C_T}{C_L}
\]
&
\[
  \eta_{\mathrm{new}}
  =
  \frac{G_T}{G_L}
  -
  \frac{C_T}{C_L}
\]
\\
Potential amplitude &
\[
  B_0^{\mathrm{old}}
  =
  \frac{\mu_L\kappa\ell_b}{\epsilon}
  \eta_{\mathrm{old}}
\]
&
\[
  B_0^{\mathrm{new}}
  =
  \frac{\mu_L\kappa\xi}{\epsilon}
  \eta_{\mathrm{new}}
\]
\\
\bottomrule
\end{array}
\]

The structural similarity is the source of the confusion.  If a reader sees only
this second table, it is tempting to think \(\ell_b\) and \(\xi\) are the same.
They are the same only if the underlying definitions also give
\[
  H_L^*=G_L,\qquad H_T^*=G_T.
\]
That is a special coefficient choice, not a consequence of the potential formula
itself.

\section{Mechanical Difference}

If one imposes the local condition
\[
  \sigma_{22}=0
\]
pointwise as the whole transverse equilibrium statement, then
\[
  C_L v-C_T\kappa Y=0,
  \qquad
  v(Y)=\beta\kappa Y,
  \qquad
  v'(Y)=\beta\kappa .
\]
Substitution into the open-circuit electric relation gives a constant electric
field,
\[
  E_2
  =
  \frac{\kappa}{\epsilon}
  \left(\mu_T-\mu_L\beta\right),
\]
and therefore the potential is only linear in \(Y\).  There is no sinh
boundary-layer term.

This is why a formula that contains a sinh term should not be interpreted as the
pure pointwise \(\sigma_{22}=0\) solution.  It already contains a higher-order
boundary-layer correction.  The issue is subtler: the older reduced expression
uses a plane-stress-like transverse closure for the bulk response and then adds a
boundary-layer correction through \(H_L^*,H_T^*\).  The new derivation instead
solves the generalized traction condition directly,
\begin{equation}
  \sigma_{22}-S_{222}'=0,
  \label{eq:generalized-traction}
\end{equation}
with the higher-order free-surface condition
\begin{equation}
  S_{222}(\pm a)=0 .
\end{equation}
Therefore the local stress \(\sigma_{22}\) is allowed to be nonzero inside the
beam.  It is balanced by \(S_{222}'\), so the generalized traction is zero.

\section{Practical Consequence}

The two formulas become exactly the same if their coefficient sets are also the
same:
\[
  \ell_b=\xi,\qquad
  \eta_{\mathrm{old}}=\eta_{\mathrm{new}},\qquad
  B_0^{\mathrm{old}}=B_0^{\mathrm{new}},
\]
and if the voltage constants are chosen consistently.  Otherwise they can look
nearly identical on paper while producing different numerical curves.

For the current plotting script, this is the main distinction:
\[
  \hbox{old comparison: coefficients based on } H_L^*,H_T^*,
\]
\[
  \hbox{new default comparison: coefficients based on } G_L,G_T.
\]
In particular, in the bare generalized-traction model, even when the mechanical
length scale \(\ell=0\), the term \(\mu_L^2/\epsilon\) can still make
\(G_L>0\), giving a finite effective boundary-layer length.

\end{document}
